\documentclass[12pt]{article}
\usepackage[utf8]{inputenc}
\usepackage[T1]{fontenc}
\usepackage{hyperref}
\title{Kambera Language Resources for K-12 Curriculum Development}
\author{Compiled by Claude Sonnet 5.5 (claude-sonnet-5-5) for Andr\'as Kornai}
\date{}
\begin{document}
\maketitle

\section{Introduction}
Kambera is an Austronesian language of eastern Sumba, Indonesia. One summary gives about 150,000 speakers, while the reference description says about 240,000. It is head-marking. It is taught as a subject in some schools. The best documentation is academic, mostly by Marian Klamer. I found no ready-made K-12 teaching materials in the searches.

\section{Primary Educational Resources}
I found no Kambera primers or textbooks in search results. The Kambera-Indonesian dictionary and the older text collections below could serve as source material for teachers.

\section{Linguistic and Reference Materials}
\begin{itemize}
\item Wikipedia, ``Kambera language''. \url{https://en.wikipedia.org/wiki/Kambera_language}
\item Marian Klamer, \emph{A Grammar of Kambera} (Mouton de Gruyter, 1998, 448 pp.): described as the first detailed, comprehensive description. Review notice: \url{https://linguistlist.org/issues/9/1342/}
\item Klamer, ``Kambera'' chapter in Adelaar and Himmelmann (eds.), 2005 (free PDF on her website). \url{https://marianklamer.org/uploads/1/2/4/7/124768088/klamer_kambera_in_adelaar_himmelmann_eds_2005.pdf}
\item Klamer 2006 chapter in Dixon and Aikhenvald (eds.). \url{https://marianklamer.lucdh.nl/uploads/1/2/4/7/124768088/klamer_2006_dixon_and_aikhenvald_eds.pdf}
\item Klamer chapter on the continuative aspect and dative clitic (Springer, 2000). \url{https://link.springer.com/chapter/10.1007/978-94-017-1580-5_3}
\item Dictionaries named in Klamer's 2005 chapter: Kapita (1982, Kambera-Indonesian) and Onvlee (1984, Kambera-Dutch). I saw no URLs for these.
\item Texts with Dutch translations from Wielenga (1913) and Onvlee (1925) are also mentioned there, without URLs.
\item A journal article on Kambera fauna vocabulary (\emph{Kulturistik}, Warmadewa University). \url{https://ejurnal.warmadewa.ac.id/index.php/kulturistik/article/view/6062}
\end{itemize}

\section{Religious and Community Resources}
I found no Kambera scripture translations or community organisations in these searches.

\section{Digital and Audio Resources}
I found no audio, video or software resources in these searches.

\section{Curriculum Development Context}
Kambera is still widely spoken, and some schools teach it as a subject. A curriculum would need new materials made with Sumbanese teachers and speakers, using Klamer's grammar and the existing dictionaries as a base. Searches in Indonesian for local-content (muatan lokal) materials would be a sensible next step.
\end{document}